% Svar på forhåndsspørsmålet, avgrensninger, policy, råd til unge, avslutning.

\begin{frame}{Vises det i statistikken? Ja, noe. Men ikke bredt.}
\begin{columns}[T,onlytextwidth]
\begin{column}{0.48\textwidth}
{\bfseries\textcolor{illInk}{Dette ser vi}}

\vspace{0.4em}
\small
\begin{itemize}\setlength{\itemsep}{5pt}
  \item De yngste faller i flere høyeksponerte yrker siden oktober 2022: utviklere $-15\,\%$, designyrker $-32\,\%$, informasjonsrådgivere $-10\,\%$ (per innbygger). De eldste i de samme yrkene vokser.
  \item Anthropic ser det samme i egne data: rundt 14\,\% lavere jobbstartrate for unge inn i eksponerte yrker (2026).
\end{itemize}
\end{column}
\begin{column}{0.48\textwidth}
{\bfseries\textcolor{illInk}{Dette ser vi ikke}}

\vspace{0.4em}
\small
\begin{itemize}\setlength{\itemsep}{5pt}
  \item Noe bredt fall i de mest eksponerte yrkene samlet, eller for 21--30-gruppen samlet.
  \item Hva som driver unge-fallet: erstattet av KI, junioroppgaver som forsvinner, strengere erfaringskrav eller seniorer som rekker mer. Der er «den perfekte stormen» av renter, konsulentkutt og normalisering etter pandemien også med.
\end{itemize}
\end{column}
\end{columns}

\vspace{0.6em}
\centering
\textcolor{illInk}{\textbf{Derfor følger vi det månedlig.}}
\end{frame}

% ---------------------------------------------------------------
\begin{frame}{Hva vi ikke sier}
\small
\begin{itemize}\setlength{\itemsep}{6pt}
  \item Indeksen er et beskrivende vekstgap, ikke en målt årsakseffekt av KI.
  \item Den dekker privat sektor; offentlig sektor inngår ikke i hovedtallet.
  \item Vi observerer utfall per yrke. Oppgaveinnholdet kan endre seg mye uten at yrkeskoden gjør det, og nye KI-roller kan være skjult i gamle koder.
  \item Yrkesstandardene er gamle: STYRK-08 er fra 2011 og bygger på ISCO-08 fra 2008, med en grunnmodell som går tilbake til ISCO-88. O*NET-oppgavene er amerikanske, vurdert i 2023. Arbeidsinnholdet i norske yrker kan avvike.
  \item Det bøter vi på slik: grove femdeler heller enn fininndeling, kryssjekk mot uavhengige mål som rangerer de samme yrkene høyt, og måling av faktiske utfall. Feil i kodingen visker ut forskjeller heller enn å skape dem.
\end{itemize}
\end{frame}

% ---------------------------------------------------------------
\begin{frame}{Politikk som treffer omstillingen}
\footnotesize
\setlength{\tabcolsep}{6pt}
\renewcommand{\arraystretch}{1.5}
\begin{tabular}{p{3.6cm} p{7.4cm} p{2.6cm}}
\textbf{\textcolor{illInk}{Sikre omskolering og inngang til jobb}} & Dagpenger under utdanning til mangelyrker, inntektssikret utdanningspermisjon, flere lønnstilskuddsplasser og stabilt statlig inntak av nyutdannede. & {\tiny\textcolor{illGray}{Kostøl \& Harang, DN 29.7.2026}} \\
\textbf{\textcolor{illInk}{Endre skatteinsentivene}} & Gjør ansettelser, opplæring og arbeidskomplementær KI mer konkurransedyktig mot ren automatisering og kjøp av maskinkapasitet. & {\tiny\textcolor{illGray}{Harang \& Kostøl, DN 29.6.2026; Acemoglu \& Johnson 2023; Brynjolfsson 2022}} \\
\textbf{\textcolor{illInk}{Finansier nye opplæringsstiger}} & Støtt sertifisert bedriftsopplæring, bransjefelles treningsløp og praksisordninger når juniorarbeidet ikke lenger finansierer læringen. & {\tiny\textcolor{illGray}{Garicano 2025}} \\
\textbf{\textcolor{illInk}{Forbered institusjonene på flere utfall}} & Følg teknologisk kapasitet, bruk og arbeidsmarkedsutfall; stresstest skatt og sikkerhetsnett og gjør tiltakene skalerbare. & {\tiny\textcolor{illGray}{Korinek 2023}} \\
\end{tabular}
\end{frame}

% ---------------------------------------------------------------
\begin{frame}{Råd til unge: skaff faglig dybde og erfaring tidlig}
\footnotesize
\setlength{\tabcolsep}{6pt}
\renewcommand{\arraystretch}{1.5}
\begin{tabular}{p{3.9cm} p{7.1cm} p{2.6cm}}
\textbf{\textcolor{illInk}{Bruk KI på virkelige oppgaver}} & Dokumenter arbeidsmåten, finn feil og kontroller kildene. & {\tiny\textcolor{illGray}{Mollick 2023; Brynjolfsson 2026}} \\
\textbf{\textcolor{illInk}{Behold selvstendig faglig arbeid}} & Gjør et eget forsøk først. Lær premisser, metode og dømmekraft. & {\tiny\textcolor{illGray}{Goldsmith-Pinkham 2026; Gans 2026}} \\
\textbf{\textcolor{illInk}{Fullfør grad eller fagbrev}} & Velg en grunnmur som åpner flere jobber, og lær KI i selve faget. & {\tiny\textcolor{illGray}{Deming 2025; Sigelman 2026}} \\
\textbf{\textcolor{illInk}{Prioriter praksis og mentorering}} & Spør hvordan arbeidsplassen trener juniorer når rutineoppgavene blir færre. & {\tiny\textcolor{illGray}{Brynjolfsson 2026; Garicano 2025}} \\
\textbf{\textcolor{illInk}{Vis hva du kan gjøre}} & Bygg noen gjennomførte arbeider og kombiner teknikk med samarbeid og kommunikasjon. & {\tiny\textcolor{illGray}{Bryan 2026; Erdil \& Besiroglu 2025}} \\
\end{tabular}
\end{frame}

% ---------------------------------------------------------------
\begin{frame}{Følg med videre}
\centering
\begin{tikzpicture}
  \useasboundingbox (0,0.2) rectangle (\textwidth,-1.3);
  \node[anchor=north west, text width=0.94\textwidth, inner sep=0pt] at (0.3,0) {%
    \small Fortrengning får mest oppmerksomhet. De to andre kreftene blir ofte oversett: at vi får gjort mer og bedre, og at nye oppgaver kommer til.};
  \fill[illRed] (0,0.05) rectangle (0.09,-1.1);
\end{tikzpicture}

\vspace{0.6em}
{\LARGE\bfseries\textcolor{illInk}{kiindeksen.no}}

\vspace{0.8em}
\large
Oppdateres hver måned når nye registerdata kommer.

\vspace{1.2em}
\footnotesize
Engelsk versjon: kiindeksen.no/en \quad $\cdot$ \quad Hernæs (Frischsenteret) og Kostøl (Handelshøyskolen BI)
\end{frame}

% ---------------------------------------------------------------
\begin{frame}{Referanser (1/2)}
\scriptsize
\begin{columns}[T,onlytextwidth]
\begin{column}{0.48\textwidth}
\begin{itemize}\setlength{\itemsep}{2pt}
  \item Acemoglu, D.\ og D.\ Autor (2011): Skills, Tasks and Technologies. \emph{Handbook of Labor Economics}.
  \item Acemoglu, D.\ og S.\ Johnson (2023): Rebalancing AI. \emph{IMF Finance \& Development}, desember 2023.
  \item Acemoglu, D.\ og P.\ Restrepo (2019): Automation and New Tasks. \emph{Journal of Economic Perspectives}.
  \item Anthropic (2026): Labor Market Impacts of AI: A New Measure and Early Evidence (Massenkoff og McCrory).
  \item Autor, D., F.\ Levy og R.\ Murnane (2003): The Skill Content of Recent Technological Change. \emph{Quarterly Journal of Economics}.
  \item Brynjolfsson, E.\ (2022): The Turing Trap. \emph{Daedalus}.
  \item Brynjolfsson, E., B.\ Chandar og R.\ Chen (2025): Canaries in the Coal Mine? Arbeidsnotat, Stanford Digital Economy Lab.
\end{itemize}
\end{column}
\begin{column}{0.48\textwidth}
\begin{itemize}\setlength{\itemsep}{2pt}
  \item Eloundou, T., S.\ Manning, P.\ Mishkin og D.\ Rock (2024): GPTs are GPTs. \emph{Science}.
  \item Google (2026): AI \& Economy ATLAS v1.0: Mapping Gemini Usage in the Economy.
  \item Hernæs, Ø.\ og A.\,R.\ Kostøl (2026): KI-indeksen. kiindeksen.no.
  \item Korinek, A.\ (2023): Scenario Planning for an AGI Future. \emph{IMF Finance \& Development}, desember 2023.
  \item Narayanan, A.\ og S.\ Kapoor (2025): AI as Normal Technology. Knight First Amendment Institute.
  \item SSB (2011): Standard for yrkesklassifisering (STYRK-08). Notater 17/2011.
  \item O*NET Resource Center, U.S.\ Department of Labor.
  \item A-ordningen via microdata.no (SSB/Sikt/NSD).
\end{itemize}
\end{column}
\end{columns}
\end{frame}

% ---------------------------------------------------------------
\begin{frame}{Referanser (2/2): kronikker, essays og intervjuer}
\scriptsize
\begin{columns}[T,onlytextwidth]
\begin{column}{0.48\textwidth}
\begin{itemize}\setlength{\itemsep}{2pt}
  \item Kostøl, A.\,R.\ og F.\,N.\ Harang (2026): KI og det stille utenforskapet. Kronikk, \emph{Dagens Næringsliv}, 29.\ juli 2026.
  \item Harang, F.\,N.\ og A.\,R.\ Kostøl (2026): Maskiner må løfte mennesker. Kronikk, \emph{Dagens Næringsliv}, 29.\ juni 2026.
  \item Garicano, L.\ (2025): AI's Becker Problem. Essay, juli 2025.
  \item Mollick, E.\ (2023): On-boarding Your AI Intern. Essay, \emph{One Useful Thing}, mai 2023.
  \item Goldsmith-Pinkham, P.\ (2026): Research in the Time of AI. Essay, mars 2026.
  \item Gans, J.\ (2026): Journopoclypse! Essay, mars 2026.
\end{itemize}
\end{column}
\begin{column}{0.48\textwidth}
\begin{itemize}\setlength{\itemsep}{2pt}
  \item Brynjolfsson, E.\ (2026): intervju, \emph{Silicon Valley Girl}, juli 2026.
  \item Deming, D.\ (2025): intervju, \emph{Managing the Future of Work} (Harvard Business School), februar 2025.
  \item Sigelman, M.\ (2026): intervju, \emph{The Ongoing Transformation}, april 2026.
  \item Bryan, K.\ (2026): intervju, \emph{Justified Posteriors}, juni 2026.
  \item Erdil, E.\ og T.\ Besiroglu (2025): intervju, \emph{Dwarkesh Podcast}, april 2025.
\end{itemize}

\vspace{0.5em}
Policy- og råd-slidene gjengir forslag fra disse kildene; de er ikke evaluert som én samlet pakke.
\end{column}
\end{columns}
\end{frame}
